\section{Origine dei target}

Per allenare il modello servono stati parziali e valori di riferimento,
chiamati \emph{target}. Per ogni stato, un procedimento separato completa
il percorso mantenendo fisse le consegne già visitate. Dal costo totale
del completamento si sottrae $g$, ottenendo il costo ancora necessario.

Sulle mappe ridotte usiamo come insegnante un branch-and-bound completo:
conserviamo il target soltanto se l'ottimo è certificato. Quando una ricerca
completa è impraticabile, usiamo invece completamenti approssimati. Si visita
ogni volta una consegna vicina e si migliora il giro con 2-opt: si invertono
tratti del percorso se la distanza diminuisce. Proviamo più partenze e
conserviamo il completamento migliore. Questo target è il costo di un giro
ammissibile, non un ottimo dimostrato.

Campioniamo percorsi parziali con molte e poche consegne residue, alternando
prefissi casuali e prefissi costruiti visitando consegne vicine. Separiamo
le mappe di allenamento da quelle di prova: stati della stessa mappa non
vengono distribuiti fra i due gruppi. La generazione dei target e il fit
sono operazioni offline, esterne al budget della ricerca valutata.
